% Standalone source-record documentation; no research attempts.
\documentclass[11pt]{article}
\usepackage[margin=25mm]{geometry}
\usepackage{fontspec,unicode-math}
\setmainfont{Latin Modern Roman}
\setmathfont{Latin Modern Math}
\usepackage{amsmath,amssymb,mathtools,xurl,hyperref}
\hypersetup{colorlinks=true,urlcolor=blue,linkcolor=blue}
\setcounter{secnumdepth}{0}
\setlength{\parindent}{0pt}
\setlength{\parskip}{5pt}
\setlength{\emergencystretch}{3em}
\newcommand{\sref}[2]{\hyperlink{source:#1:#2}{[S#2]}}
\newcommand{\cataloguescope}[1]{\paragraph{Source record.}#1}
\newcommand{\floor}[1]{\left\lfloor#1\right\rfloor}
\begin{document}
% UnsolvedMath ID 20001186; source catalogue code AIM-COMPUTATION-0024
\setcounter{section}{20001185}
\section{Exact endpoint and stability census for the symmetric triangle}\label{20001186}
{\small\sffamily UnsolvedMath ID 20001186 · Computer Science}
\subsection{Definitions and mathematical statement}

\paragraph{Shared notation.}Standard number systems are $\mathbb N,\mathbb Z,\mathbb Q,\mathbb R,\mathbb C$. Specialized notation and index ranges follow the cited source and the definitions below.


For an integer matrix $A$ and statistic vector $b$, the nonnegative integer fiber is $\mathcal F_b=\{x\in\mathbb N^d:Ax=b\}$. A Markov basis is a finite set of integer moves in $\ker A$ connecting every such fiber by moves that remain nonnegative. Its binomial generators form the toric ideal $I_A=\langle x^u-x^v:Au=Av\rangle$. Network tomography infers internal network characteristics from indirect endpoint measurements. A rational generating function compactly encodes integer points through monomials. Seemingly unrelated regression uses several regression equations with correlated errors and can have a nonconcave likelihood.
\paragraph{Record scope and qualifications.}
The catalogue title is retained as metadata. Named certificates, contracts, bounds or frameworks in that title are not asserted as results of the original source extract.
The actual blocking equations and meaning of statistically interesting solutions are absent. The later questions have different objects and models; no endpoint or stability census is reconstructed. The supplied catalogue marks this record partially solved; its present research status was not independently checked.
The source text below is rendered literally where necessary; damaged formulas and composite extracts are not silently repaired into a new mathematical claim.
\paragraph{Statement recorded in the supplied source.}
Question 10. Are there always only 0, 1, or 2 statistically interesting solutions to the blocking probability equations in network tomography?

Ruriko Yoshida

Ruriko Yoshida discussed short encodings of the set of points in polytopes by means of rational generating functions, and the applications of rational generating functions for giving short encodings of Markov bases. Ian Dinwoodie asked, "Is there a way to extract (random) monomials from the gener-ating functions?" Akimichi Takemura asked, "Is there a way to extract only the absolutely essential Markov basis elements that are needed to connect a given fiber of the form \{Ax = b|x ∈ Nd\}?Ruriko answered,"We are still attempting to solve these problems." Henry Wynn asked, " Can these techniques be used to approximate the number of integral points inside an arbitrary convex set?"

Mathias Drton

Mathias Drton discussed multimodality of the likelihood function for the model of seemingly unrelated regression (SUR). These can be thought of as conditional independence models induced by the chain graph with two sets of nodes, ( X1,..., X n) and ( Y1,..., Y n)and directed edges from Xi → Yi and undirected edges between any pair Yi − Yj.6 \sref{20001186}{2}
\paragraph{Source status.}
The supplied catalogue labels this source record partially solved; that metadata was checked against the supplied snapshot on 2026-10-11. The label is preserved for provenance and is not an independent verification of present open status or mathematical correctness.
\subsection{Short English statement}
The excerpt combines blocking-probability equations, generating-function questions, and regression likelihood discussion.
\subsection{Sources}
\begin{itemize}
\item[\hypertarget{source:20001186:1}{S1}] Ulam AI and the UnsolvedMath Contributors, UnsolvedMath: A Curated Collection of Open Mathematics Problems, record 20001186 (AIM-COMPUTATION-0024). Catalogue curation and metadata: CC BY 4.0. The source extract is retained; definitions, English summaries and qualifications are editorial additions.
\url{https://huggingface.co/datasets/ulamai/UnsolvedMath}.
\item[\hypertarget{source:20001186:2}{S2}] AIM Workshop Problem Lists, source linked by catalogue record AIM-COMPUTATION-0024.
\url{https://aimath.org/WWN/compalgstat/compalgstat.pdf}.
\end{itemize}
\paragraph{Preparation and review.}
The original source and UnsolvedMath contributors retain attribution for the source extract. Codex (GPT-6; exact serving revision not exposed) assisted with definitions, English summaries, scope notes and formatting. No new human mathematical authorship or independently verified solution is asserted. Permanent source ID 20001186 and code AIM-COMPUTATION-0024 checked against the complete supplied registry and current repository. Distinct upstream records, including repeated questions and broad or damaged extracts, are retained for complete coverage of the source CS category. A repeated record is not a new independent conjecture. Complete mathematical-equivalence and current-literature review has not been performed. Preparation and final review follow contribution prompts C and D, with source gaps disclosed. Structural and rendering checks do not establish a new theorem or open status.
\label{end:20001186}
\end{document}
